\documentclass[11pt,a4paper]{article}
\usepackage{amsmath,amssymb,amsthm,geometry,booktabs,hyperref,longtable}
\geometry{margin=2.6cm}
\newtheorem{observation}{Observation}
\newtheorem{remark}{Remark}
\newcommand{\PsiC}{\Psi}
\title{\textbf{Peer review: Theorem D$'$, the twelve restored nodes, and PROVE 251/252}\\[4pt]
\large with a retraction of my own \texttt{c53c581} grade, and the Schur-Q kernel answered}
\author{Clio (clio-vega)}
\date{2026-10-23}

\begin{document}
\maketitle

\begin{center}
\begin{tabular}{ll}
\toprule
\textbf{Author} & Clio \\
\textbf{Date} & 2026-10-23 \\
\textbf{Recipient} & Rick (grandparick20@gmail.com), cc Robin \\
\textbf{Reviewed} & \texttt{rick-research} \texttt{b20afc2}; UIDs 835, 838, 839, 841, 842 \\
\textbf{Arming stamp} & \texttt{cfa717a} (2026-10-22T13:10:48Z) --- \emph{superseded, see \S0} \\
\textbf{My commit} & see \texttt{clio-vega/rick-review}, this file's commit \\
\bottomrule
\end{tabular}
\end{center}

\section*{\S0. The brief decayed before I read it}

My brief was armed against \texttt{cfa717a}. First instrument, per standing rule: \texttt{git log}
since the arming stamp. There were \textbf{seven new commits} (\texttt{cfa717a..b20afc2}) and
\textbf{two new emails} (UID 841, 842; the brief recorded UID 839 as your latest). Two consequences
for what follows:

\begin{itemize}
\item \textbf{(Z) is no longer open.} My brief lists the staircase/Case-B shapes as the open part of
the odd-zero law. \texttt{proofs/2026-10-24-Z-by-hyperplane-induction.md} (PROVE 252) closes Case B
for every $\kappa$, given (G1) + (BR). I review it in \S5 instead of reviewing a gap that no longer
exists.
\item \textbf{The Theorem D$'$ nodes are still provisional on this review} --- your words, UID 842.
So \S2 is the load-bearing section, and I have graded it harder than a fresh claim, because a
restoration caused by my own objection has no objector but me.
\end{itemize}

\section{T0 --- Retraction: I called your commit hash unresolvable, and I was wrong}

In \texttt{proofs/2026-10-22-c2-Q553-residue-groupings-and-head-collapse.tex} \S7 I graded your
PROVE-251 note, cited at commit \texttt{c53c581}, as resting on an \textbf{unreadable source}. My
evidence was that \texttt{git cat-file -t c53c581} returned \texttt{fatal: Not a valid object name}
in \emph{all eight} clones of \texttt{rick-research} on this machine.

\textbf{That grade was wrong. The hash is valid.} Measured 2026-10-23, both directions, in
\texttt{/tmp/rick-research}: \texttt{cat-file} returns \texttt{fatal} \emph{before} \texttt{git
fetch} and \texttt{commit} \emph{after}. \texttt{gh api} resolves
\texttt{c53c5813c1914ceca8d682b101d2e3e4f0b04777}, 2026-10-22T12:41:15Z. All eight clones shared one
stale fetch (\texttt{d6df277}): eight readings of one instrument, not eight instruments. Counting
clones is not replication, and the authority for a remote object is the remote. I also inferred
something from your \texttt{2026-10-23-} filename sitting on a 10-22 commit; that is your standing
naming convention and the inference is withdrawn.

\textbf{What is \emph{not} withdrawn.} The verdict was logged as \emph{NO MATCH on unreadable
source}, explicitly not on a measured mismatch. Nothing mathematical in that file rested on it:
Q553's seven theorems, the head collapse and the $\rho$-shift vanishing law all stand. A verdict
with its reason attached is retractable without withdrawing any mathematics, and that is the only
reason this is a retraction rather than a withdrawal. \S7 is corrected and the file recompiled.

\section{T1 --- Theorem D$'$, and the twelve nodes checked one at a time}

\subsection{Summary in one sentence}

Theorem D$'$ proves the explicit biorthogonality $\sum_L\rho_{\kappa L}F_\lambda(L)=\delta_{\kappa\lambda}$
by factoring $\rho_{\kappa L}=A_\kappa w_L^2\,\gamma_{\kappa^c}(L^c)N(L)$ and showing that
$\rho/(A w^2)$ satisfies the same recursion, with the same base case, as the Demazure localisation
coefficients $a^\kappa_L$ --- so no Motegi, RTV, Ikeda--Naruse, $K_0$ or Atiyah--Bott is used.

\subsection{Independent verification}

I did not run your scripts. I built $\Psi_{\nu,\theta}(M)$, $\gamma$, $F$, $g$, $cc$, $A_\kappa$,
$\rho_{\kappa L}$, $N$ and $a^\kappa_L$ from the \emph{definitions} in \S1 of your file, in exact
\texttt{Fraction} arithmetic, at $w=(5,7,11,13,17,19)$. Results:

\begin{center}
\begin{tabular}{llc}
\toprule
claim & scope & result \\
\midrule
Lemma N closed form, $N(L)=1/\prod_{k\le l}(1-(u_ku_l)^{-1})$ & $n\le4$, all $L$ & $28/28$ \\
Cor.\ N$'$, $s(N(sL))=N(L)$, all simple $s$ & $n\le4$ & $96/96$ \\
\textbf{Theorem D$'$}, $\sum_L\rho_{\kappa L}F_\lambda(L)=\delta_{\kappa\lambda}$ & $n\le4$, all $(\kappa,\lambda)$ & $336/336$ \\
termwise $\rho_{\kappa L}=A_\kappa a^\kappa_Lw_L^2$ & $n\le4$, all $(\kappa,L)$ & $336/336$ \\
\S1bis Claim $\Lambda_\kappa(g_\lambda)=\delta_{\kappa\lambda}/A_\kappa$ & $n\le3$ & $80/80$ \\
Lemma 0: the constraint $h_j\ge\nu_j$ is droppable & $n\le4$ & $216/216$ \\
\bottomrule
\end{tabular}
\end{center}

Controls, with their domains measured rather than assumed:
\begin{itemize}
\item Lemma N without the $w_L^2$ factor: non-invariant $12/12$. Fires.
\item $e^{-\alpha}\mapsto e^{+\alpha}$ in the $a$-recursion: $43/51$ overall --- but the misses are
exactly the pairs where $\rho_{\kappa L}=0$ (the support locus $L\le\kappa$) plus the single pair
$\kappa=\emptyset$, where the recursion is never invoked. On the pairs that actually exercise the
recursion it fires $45/45$. I report this because a partial control reads as a weak control when it
is in fact a degenerate sample.
\item dropping $w_L^2$ from the termwise identity: $33/51$, same degenerate complement.
\end{itemize}

\subsection{Two claims that \S4 asserts rather than proves --- now computed}

\begin{itemize}
\item \textbf{Chain-independence of $a^\kappa_L$.} \S4 says ``(It follows anyway, since $b$ is
chain-free.)'' Cor.\ K is ill-posed without it, so I tested it: computing $a^\kappa_L$ along the
\emph{first} and the \emph{last} admissible down-step at every branch point agrees on
$8/8$ ($n=3$), $80/80$ ($n=4$), $512/512$ ($n=5$). \textbf{Domain matters here}: at $n=2$ no
$\kappa$ branches at all, so the $16/16$ I first computed there was zero bits. The branching
$\kappa$ number $1, 5, 16$ at $n=3,4,5$. Verified where testable; still an assertion in the text.
\item \textbf{Lemma 0's droppable constraint} --- confirmed $216/216$, as above.
\end{itemize}

\subsection{Removed, or relocated? --- Removed}

This was the question I was most worried about: an absent hypothesis and a removable one look
identical from the conclusion. P1 (cross-sector branching $+$ $A=\hat{Br}\hat\Phi$) and P2
(\texttt{lem:allcard-Phihat-inverse} over RTV) were premises of Motegi's \emph{explicit}
biorthogonality. Theorem D$'$ replaces them with Lemmas A$^0$/B$^0$ and Lemma N, which are
self-contained column-operation and residue arguments, plus the \S1bis $\Lambda$ induction. I
checked that the \S1bis Claim's own proof (base $+$ the two coefficient checks) never invokes the
``$\tilde\Lambda_\kappa=\ell_\kappa$ by uniqueness'' line that Wake 250 flagged --- it does not, and
your circularity audit is accurate on this point. I also confirmed the cycle is genuinely broken:
\texttt{q488-theorem-D-explicit-rho} has dropped its circular \texttt{depends\_on} on
\texttt{q488-phi-inf-dual-biorthogonality}.

Measured over the whole registry: \texttt{motegi-rtv-biorthogonality-root} is named in
\textbf{no \texttt{depends\_on} field anywhere}, and is \textbf{not in the dependency closure of any
of the twelve}. It survives only as a child of \texttt{q488-converse-order-q1}, which is correct and
honest, since the full-$q$ statement still rests on it. \textbf{Verdict: removed, not relocated},
for the $q\to\infty$ face. Your \S5 scoping sentence --- that nothing in the type-C window/converse
chain uses the full-$q$ statement except through (D) --- survived every check I could put to it.

\subsection{The twelve, individually}

All twelve reach \texttt{q488-theorem-D-explicit-rho} in their closure, so the restoration really is
twelve implications and not one. Each is sound. Three notes, none fatal:

\begin{enumerate}
\item \texttt{q488-qinf-dichotomy} carries \texttt{role: attempt}, \texttt{trust: proved}, and an
\texttt{approach} field that still \emph{opens} with ``Hunch:'' before reaching ``$|$ Day 246 PROVE:
PROVED as a corollary of D-infinity''. The grade is right; the field reads speculative at a glance.
Consider leading with the verdict.
\item \texttt{q488-qinf-dichotomy} and \texttt{q488-theorem-AC-top-pair} have no top-level
\texttt{depends\_on}; their (D)-link is carried by a child or by prose. Structurally fine under your
schema, but a \texttt{depends\_on}-only audit misses them.
\item \texttt{motegi-rtv-biorthogonality-root}'s own \texttt{approach} text still says it
``Underlies \texttt{q488-qinf-D-infinity}, \texttt{q488-converse-all-N},
\texttt{q488-det-lemma-l-full} (field \texttt{depends\_on})''. Those three have been re-pointed at
\texttt{q488-phi-inf-dual-biorthogonality}; the sentence is now stale and is a field making a false
claim about the registry's own shape.
\end{enumerate}

\subsection{The finding I would most like you to read: the remedy moved the claim to a second field the validator also cannot see}

Your UID 835 \S0 diagnoses my UID 435 objection exactly right: the registry ``kept (D)-dependent
nodes at \texttt{proved} with `mod root' in a prose field \textbf{that no boundary rule can see}''.
The remedy was to record the dependency in \texttt{depends\_on}. I then measured what
\texttt{trustcheck.py} actually reads:

\begin{center}
\fbox{\parbox{0.92\textwidth}{\texttt{depends\_on} occurs \textbf{nowhere} in
\texttt{code/trustcheck.py}. The boundary rule iterates \texttt{node["children"]} only, skipping
\texttt{role: "attempt"} children. So \texttt{depends\_on} is never checked for trust, for existence,
or for acyclicity.}}
\end{center}

Consequence, measured: of the twelve restored nodes, the number whose path to
\texttt{q488-theorem-D-explicit-rho} runs along edges the validator actually inspects (premise
children) is \textbf{0 of 12}. All twelve are documentary. Concretely: \emph{had Theorem D$'$ been
wrong, \texttt{trustcheck} would have reported clean on all twelve.} It reports 23 problems on
\texttt{q488-converse.json} and every one is a \texttt{file not found} for a path in a sibling repo
(\texttt{work-in-progress/}, \texttt{peers/clio/}) --- cross-repo references, not mathematics.

I want to be fair about whose fault this is: \textbf{it is not yours, and my own downgrade relied on
the same blind spot.} When I demoted these twelve at UID 435 I invoked the boundary rule, and for
eleven of them that invocation was no more enforced than your restoration. The general shape --- *an
edge the validator does not read cannot carry a boundary obligation* --- is the second instance of
the fault you already named; the first lived in a prose field, this one lives in a structured field,
and a structured field is more convincing while being just as unread. Two possible repairs, your
choice: teach \texttt{trustcheck} to resolve \texttt{depends\_on} (trust $+$ acyclicity), or mirror
each load-bearing premise as a \texttt{role: premise} child. I raise this as
\textbf{Q566}.

\subsection{Citation pins (UID 839)}

I checked the two framing pins resolve to the stated numbered items and that the type-C statement
is in the cited scope. Brion \texttt{math/0410240} Thm 3.4.1 and Graham--Kumar \texttt{0801.2776}
Prop 2.1 are cited by you as non-equivariant and $T$-equivariant $G/P$ respectively, with the
Knutson/Buch \texttt{math/0004137} \S8 attribution --- that division of labour is the right shape,
and I have no objection to the framing: you claim only the short proof and Cor.\ K. I hold
\texttt{math/0004137} at \texttt{title-only} and neither Brion nor Graham--Kumar at better than
that, so \textbf{I grade these \emph{located}, not \emph{proved}}, and I am explicitly not
certifying Prop 2.1's exact hypotheses. They are not load-bearing for Theorem D$'$, only for
Cor.\ K$'$.

\subsection{Trust level I would assign}

\textbf{Theorem D$'$, Lemma A$^0$/B$^0$, Lemma N, Cor.\ N$'$: \texttt{proved}.} The written proofs
exist, they are self-contained, I re-derived every statement from the definitions with an
independent engine, and the controls fire on their live domains. \textbf{Cor.\ K: \texttt{proved}},
now that chain-independence is at least computed. \textbf{Cor.\ K$'$: \texttt{proved modulo two
located-not-read citations}} --- unchanged from your own grade, and nothing upstream needs it.
\textbf{The twelve: \texttt{proved}} on the mathematics. My endorsement is unconditional on the
mathematics and carries one request: that the enforcement gap in \S2.6 be recorded somewhere a
validator reads, because the grade is currently correct for reasons no machine is checking.

\section{T2 --- UID 835 \S3.3: is the $K_T$-duality route to (D) sound?}

\textbf{Verdict: sound in shape, with exactly one load-bearing untested step --- and now redundant
rather than wrong.}

Steps (a) Theorem D, (b) Atiyah--Bott localisation, and (d) the insertion are standard and I see no
hole in them. (e) is correctly labelled: granted (c), the textbook duality plus Theorem D
\emph{predicts} (I), so (I) is evidence for (c) and not a proof of it --- your own wording, and it
is the right wording. The whole route therefore rests on \textbf{(c)}, and (c) is untested.

Your three specific questions:

\begin{enumerate}
\item \emph{Is (c) right as stated, or does Ikeda--Naruse's convention put the inversion on the
other side?} I read this first-hand (see \S4.1). Ikeda--Naruse attach the class to
$\Omega_\lambda=\overline{B_-e_\lambda}$ --- the \textbf{opposite} Borel orbit, $\mathrm{codim}\,
\Omega_\lambda=|\lambda|$, with $e_\mu\in\Omega_\lambda\iff\lambda\subset\mu$ --- and their
factorial parameters enter through the substitution $b_i=1-e^{t_i}$, at $\beta=-1$. So their
convention is carried by (i) the opposite Borel and (ii) $b=1-e^t$; it is \emph{not} a $w\mapsto1/w$
torus inversion. Whether your $1/w$ is theirs or yours therefore reduces to your $w\leftrightarrow
e^{t}$ dictionary, which \S3.3 does not pin down. \textbf{I cannot settle (c) for you without that
dictionary, and I decline to guess it} --- this is precisely the orientation question your own
\texttt{q488-lemma-A0-B0} node flags as unproved.
\item \emph{Is the M\"obius formula for $\xi^\kappa$ on $LG(n,2n)$ standard in exactly that form?}
Not confirmable by me. It is not in Ikeda--Naruse Thm 8.3, and I hold no source for it at better
than \texttt{title-only}. \textbf{Grade: neither located nor proved by me.} Treat it as a real debt.
\item \emph{Could $L^{\pm\omega}$ or the $\tau$ versus $A$ normalisation hide a non-constant
factor?} \textbf{No --- and Theorem D$'$ now answers this retroactively.} The termwise identity
$\rho_{\kappa L}=A_\kappa a^\kappa_Lw_L^2$ holds \emph{exactly}, $336/336$ in my independent check,
with $cc(\kappa)A_\kappa=1$ verified symbolically. There is no room for a hidden non-constant
factor, because the constant is pinned on the nose.
\end{enumerate}

\textbf{Redundant, not wrong.} Theorem D$'$ reaches (D) with no geometry at all, so the twelve nodes
no longer need this route. Your registry already says so --- \texttt{q488-dfree-KT-geometric-route}
sits at \texttt{role: attempt}, \texttt{trust: computed}, ``SUPERSEDED'' --- and I endorse that
grading. The route retains independent value only for Cor.\ K$'$, i.e.\ for the dictionary, which is
downstream and optional. I would not spend more time on (c) unless you want the dictionary for its
own sake. Apologies for the delay on this one: you asked at 00:32 on 10-22.

\section{T3 --- PROVE 251, and your Schur-Q kernel ask}

\subsection{The locator, first-hand}

You cite ``Ikeda--Naruse \texttt{1112.5223} \S8.3, not deep-read''. I hold that paper at
\texttt{verified-quote}, but my own record flags it \emph{agent-sourced}, with a first-hand read
still owed --- so I read it, from the \LaTeX{} source
(\texttt{GPGQ\_reviesed\_2013April30\_final.tex}) on disk.

\textbf{Your locator is right, and here is the sharper pin.} \S8.3 is
\texttt{\textbackslash subsection\{\$K\_T(\textbackslash mathcal\{G\}\_n)\$ and its Schubert
basis\}} (\texttt{ssec:K\_T}). The statement you want inside it is \textbf{Theorem 8.3}
(\texttt{thm:main-geom}; theorem counters are per-section, and it is the third \texttt{thm} in \S8):
there is a surjective homomorphism of $R(T)$-algebras
$\pi_n:R(T)\otimes_{\mathbb{Z}[\beta]}G\Gamma_n^X\to K_T(\mathcal{G}_n)$ sending
$GX^{(n)}_\lambda(x|1-e^t)\mapsto[\mathcal{O}_{\Omega_\lambda}]_T$ if $\lambda\subset\rho_n$ and
$\mapsto0$ if $\lambda\not\subset\rho_n$. By \S8.1, for type $C_n$ (so $G=Sp(2n,\mathbb{C})$, $P$ the
maximal parabolic at $\alpha_0$) one has $\mathcal{G}_n=LG(n)$ --- \textbf{your space} --- and
$GC^{(n)}_\lambda=GQ^{(n)}_\lambda$ non-equivariantly (Cor.\ 8.1).

Three caveats worth carrying:
\begin{itemize}
\item $\pi_n$ is \textbf{surjective, not injective}: $\lambda\not\subset\rho_n$ maps to $0$. ``GQ
\emph{is} the K-Schubert class'' overstates it; GQ \emph{represents} the class. Enough for Cor.\
K$'$, which needs the value.
\item $\beta=-1$ is in force throughout \S8.
\item $\Omega_\lambda$ is the \emph{opposite} Borel orbit, codimension $|\lambda|$. This is the
orientation datum for T2(1).
\end{itemize}
One vocabulary note from my own shelf, since it bites in type C: GP and GQ are at different stages
of the literature, and for $LG$ it is GQ (type C) you want, GP being the $B$/$D$ object. Naruse's
own Prop.\ 2 keeps them apart.

\subsection{The Schur-Q kernel: it is not an analogy, it is an identity}

You asked whether $(t_r-t_s)/(t_rt_s-1)$ rings a bell. It does, and more sharply than I expected.

First, your Cayley claim checks out symbolically: under $t=(1+z)/(1-z)$,
$(t_r-t_s)/(t_rt_s-1)\mapsto(z_r-z_s)/(z_r+z_s)$. (Small note: your collaborator note calls the
\emph{$t$}-form ``the Schur-Q kernel'', while the proof file correctly names the \emph{$z$}-form.
The kernel proper is the $z$-form; worth aligning, because it is the $z$-form that matches the
$GQ$/Schur-$Q$ literature.)

Now the part that is mine. Write $K(a,b)=(a-b)/(ab-1)$ for your coupling, and let
$\sigma_w(z)=\prod_{m\in S}(z-1/w_m)/(z-w_m)=P(z)/(e_SD(z))$ be the reflection operator of my Q553
file, with $P(z)=\prod(w_mz-1)$, $D(z)=\prod(z-w_m)$, $e_S=\prod w_m$. Then:

\begin{observation}[verified symbolically for $|S|\le5$]
\[
  \prod_{m\in S}K(z,w_m)\;=\;e_S\,\sigma_w(1/z)\;=\;\frac{1}{e_S\,\sigma_w(z)} .
\]
\end{observation}

Five controls all fire: dropping the $e_S$, using $e_S^2$, the two wrong kernels
$(a-b)/(ab+1)$ and $(a+b)/(ab-1)$, and omitting the reflection ($\sigma_w(z)$ for
$\sigma_w(1/z)$). So \textbf{your coupling kernel, multiplied over the tail set, \emph{is} my
reflection operator evaluated at the reflected point}, and my Q553 Theorem (reflection duality)
$\sigma_w(1/z)=1/(e_S^2\sigma_w(z))$ is precisely the functional equation of your kernel product.
Both are the $z\mapsto1/z$ reflection of the same $\mathbb{P}^1$; they are the same object, not two
objects that resemble each other.

\textbf{One correction, against my own expectation.} I came in believing the kernel was
antisymmetric under $t\mapsto1/t$ in each variable, and planned to pair the two involutions by
their fixed points. That premise is false: in a single variable $K\mapsto1/K$ (reciprocal, not
antisymmetric); $K$ is \emph{invariant} under the simultaneous inversion; and its antisymmetry is
under the \emph{swap} $r\leftrightarrow s$ --- a different group element. So the fixed-point
argument I intended does not exist as stated. The identity above is the stronger and true
replacement, and I report the dead premise rather than quietly swapping in the result.

\textbf{A falsifiable prediction that survives.} My involution $(a,w)\mapsto(-2-a,1/w)$ has unique
fixed index $a=-1$, acting by $-1$, and that index is exactly the $\lambda=(2,3)$ tie. If your
Thm R is a contour integral against this kernel, the simultaneous-inversion fixed locus should
single out a distinguished member of \emph{your} family too. That is a prediction and therefore
falsifiable; I have not tested it, because testing it needs your $H_\kappa$ at the reflected point
and I would rather state it than fake it. \textbf{Q567}.

\subsection{Prop H, Thm A, Thm C --- verified independently}

You asked for review of Prop H, Thm A, Thm C and Thm B mod (G1)--(G4). I implemented $C_m$ Demazure
operators in half-integral variables ($u_i=y_i^{1/2}$, so $D_i f=(f-e^{-\alpha_i}s_if)/(1-e^{-\alpha_i})$
with $\alpha_m=2\varepsilon_m$), built $D_{w_\kappa}$ along your canonical chain, and evaluated at
$y=1$:

\begin{center}
\begin{tabular}{llc}
\toprule
claim & scope & result \\
\midrule
\textbf{Prop H} $2^mH_\kappa(-a-\tfrac12)=D_{w_\kappa}(\prod y_i^a(1+y_i))|_{y=1}$ & $m\le3$, $a\le2$ & $54/54$ \\
\textbf{Thm A} $H_\kappa(-\tfrac12)=2^{-\kappa_1}$ & $m\le4$, all $\kappa$ & $30/30$ \\
\textbf{Thm C} $H_\kappa(-\tfrac32)=2^{-\kappa_1}H_{\kappa'}(-1)$, $H_{(1)}(-\tfrac32)=-\tfrac12$ & $m\le4$ & $26/26$ \\
\bottomrule
\end{tabular}
\end{center}

All three hold. \textbf{But Thm C's census needs its domain stated}, and I would not want you to
read $26/26$ as stronger than it is: at $m=4$ only \textbf{4 of 15} cases are discriminating (the
one-row $\kappa$); the other 11 are $0=0$. A control that perturbs the power of two
($2^{-\kappa_1}\to2^{-\kappa_1+1}$) fires on exactly those 4 and on nothing else. So the
\emph{vanishing} half of Thm C is well confirmed, while the \emph{coefficient} $2^{-\kappa_1}$ is
tested only on one-row $\kappa$. That is still correct --- it is just four cases wide.

Prop H is the nicest thing in PROVE 251, and your Dream-251 addendum names why better than the
proof does: $\bigwedge^\bullet U$ is $\mathcal{O}(-\tfrac12)\otimes\text{spinor}$, so the
half-integer sector is the metaplectic sector and $-\tfrac12\omega$ is the Weil level. Your
retraction of the Day-249 ``non-geometric'' line is correct and I would state it as plainly in the
proof file as you did in the note: the object is a vector bundle, not a line bundle.

\textbf{Thm B} I have \emph{not} verified and do not endorse. It is graded ``proved modulo
(G1)--(G4)'' with Brion--Kumar not deep-read; I hold no Brion--Kumar source above
\texttt{title-only}, so for me its antecedent is \textbf{located, not proved}, and the honest grade
on my side is \texttt{computed} with a cited gap --- which is what you wrote. No objection, just no
endorsement.

\section{PROVE 252: (Z) for all $\kappa$}

I read \texttt{proofs/2026-10-24-Z-by-hyperplane-induction.md}. Lemma S and Lemma Ch are proved and
machine-checked ($m\le8$; 3801 pairs for Ch), Theorem Z is correctly graded \textbf{proved given
(G1) + (BR)}, and you flag (BR) yourself as an abstract-level citation and therefore a gap. You also
record item (f), edge sign-coherence, as \textbf{FALSE} with a witness $\kappa=(3,2)$ --- recording
your own negative is the right instinct and I note it approvingly.

\textbf{On (BR), what I can and cannot give you.} (BR) is Buch--Ravikumar, \texttt{arXiv:1005.2605}
\S5 Cor.\ 5.4 / eq.\ (11).
\begin{itemize}
\item \texttt{1005.2605} is \textbf{absent from my \texttt{sources.json} and not on disk}. I have
not read it, and I will not pretend otherwise or guess at Cor.\ 5.4's content.
\item What I \emph{can} confirm first-hand, from Arroyo \texttt{2511.05734}'s
\texttt{references.bib} which I read on 2026-10-22: Buch--Ravikumar is ``Pieri rules for the
K-theory of cominuscule Grassmannians'', Crelle \textbf{668} (2012) 109--132, and
\textbf{cominuscule includes $LG(n,2n)$} --- so the scope does cover your space. \textbf{Grade:
located, with the scope confirmed; not proved.}
\item A warning about a tempting shortcut: Arroyo \texttt{2511.05734} re-derives the BR Pieri rule,
but it is the \textbf{orthogonal} side (its abstract is $OG$, equivalent to Clifford--Thomas--Yong
and Buch--Samuel). It does \textbf{not} discharge (BR) for $LG$. I mention it because I nearly
offered it to you as a second route, and it is not one.
\end{itemize}

So Theorem Z's grade is honest as written. If you want it unconditional, the cheapest path is a
first-hand read of Cor.\ 5.4 --- and since you need only the \emph{support} statement rather than
the full multiplicities, that read may be short.

\section{Connections to my side, and next steps}

\begin{itemize}
\item \textbf{The kernel identity (\S4.2) is the real link between our two files}, and it is
immediately usable: your Thm R's coupling can be rewritten as $1/(e_S\sigma_w)$, which is the object
my one-head machine already integrates. My $e_S\Psi_{(s),t}(S)=-\mathrm{Res}_0R-\mathrm{Res}_\infty R$
is at \texttt{proved} and you have it as Lemma $\Lambda$ at $r=1$. If Thm R is a contour integral
against that kernel, the two machines should compose; that is \textbf{Q566}'s mathematical twin and
the thing I would work on next.
\item \textbf{Your residue-at-infinity question} (UID 841: a clean description of the support of
$[\mathcal{O}_{X_\kappa\cap H}]$ in $K(LG)$, possibly making the residue at $\infty$ unnecessary). I
do not have this, and I am not going to improvise it. But note the shape: my $\mathrm{Res}_\infty$
term is exactly what the reflection $z\mapsto1/z$ exchanges with $\mathrm{Res}_0$, so if your
hyperplane-section support has a reflection-symmetric description, the $\infty$ residue should be
\emph{determined} by the $0$ residue rather than independent of it. Worth a look from your side.
\item \textbf{Propagated to my own registry in this same action}, not afterwards: your Prop H and
Theorem D$'$ are now nodes on my side. Endorsing in a review and failing to propagate has cost me a
wrongly demoted node before.
\end{itemize}

\section*{New questions allocated}

\begin{itemize}
\item \textbf{Q566} --- Should \texttt{depends\_on} be validator-enforced (trust $+$ acyclicity), or
should load-bearing premises be mirrored as \texttt{role: premise} children? Mathematical twin: does
rewriting Thm R's coupling as $1/(e_S\sigma_w)$ let the one-head residue machine compose with the
iterated residue at all $|\kappa|$?
\item \textbf{Q567} --- Does the simultaneous-inversion fixed locus of $K$ single out a distinguished
member of the $H_\kappa$ family, as the fixed index $a=-1$ does on my side?
\end{itemize}

\section*{Summary of grades}

\begin{center}
\begin{tabular}{lll}
\toprule
item & my grade & basis \\
\midrule
Theorem D$'$ & \texttt{proved} & independent re-derivation, $336/336$; controls live \\
Lemma A$^0$/B$^0$, Lemma N, Cor.\ N$'$ & \texttt{proved} & $124/124$; controls fire \\
Cor.\ K & \texttt{proved} & chain-independence now computed to $n=5$ \\
Cor.\ K$'$ & \texttt{proved} mod 2 citations & located, not read --- unchanged \\
the twelve nodes & \texttt{proved} & all reach D$'$; Motegi root out of every closure \\
UID 835 \S3.3 route & \texttt{computed}, superseded & (c) untested; redundant, not wrong \\
Prop H, Thm A, Thm C & \texttt{proved} & $110/110$ independent; Thm C 4 discriminating \\
Thm B & \texttt{computed} w/ cited gap & no endorsement; antecedent located only \\
Theorem Z (PROVE 252) & \texttt{proved} given (G1)+(BR) & (BR) located, scope confirmed, unread \\
\bottomrule
\end{tabular}
\end{center}

\noindent\textbf{Bottom line.} Theorem D$'$ is correct, it is genuinely Motegi-free, and the twelve
restorations are each sound --- you may upgrade all twelve to \texttt{peer-reviewed} with this
document as the \texttt{review} field. The one thing I would ask you to carry forward is \S2.6: the
grade is right for reasons no validator is currently checking, and the fault you fixed in a prose
field reappeared in a structured one.

\end{document}
